\begin{frame}{Методы Рунге--Кутты}
	\tiny
	\begin{itemize}
		\item Решение задачи Коши
		\[
		y' = f(t,y), \qquad y(t_0)=y_0.
		\]
		
		\vspace{0.2cm}
		
		\item Общий вид $s$-стадийного метода:
		\[
		y_{n+1}=y_n+h\sum_{i=1}^{s} b_i k_i.
		\]
		
		\vspace{0.2cm}
		
		\item Коэффициенты метода задаются таблицей Бутчера
		\[
		\begin{array}{c|c}
			c & A \\
			\hline
			& b^T
		\end{array}
		\]
		
		\vspace{0.2cm}
		
		\item Для методов высокого порядка возникает система нелинейных условий порядка.
		
		\vspace{0.2cm}
		
		\item После выполнения условий порядка остаются свободные параметры:
		\[
		\theta=(c_2,c_4,c_5,c_7).
		\]
		
		\item Задача работы:
		подбор параметров $\theta$ для осциллирующих задач.
	\end{itemize}
	
\end{frame}